%\documentclass[10pt,respuestas,a4]{examenEPN}
\documentclass[10pt,a4]{examenEPN}

\usepackage{comandosEPN}
\usepackage{esint}
\usepackage{tikz}
\usepackage{systeme}
\usetikzlibrary{automata, positioning}
\usepackage{multicol}
\usepackage{enumitem} 
\usepackage{multicol}

\usepackage{pstricks,pst-plot}

\materia{Álgebra Lineal}
\examen{Hoja de ejercicios Valores y Vectores Propios}

\autor{Departamento de Formación Básica}
\fecha{Semestre 2024-B}
\logouno[1.6cm]{Logo_EPN.eps}
\logodos[2.6cm]{Logo_DFB.eps}

\begin{document}

\encabezado
%\notapie{Ejercicios sugeridos: 
%Valores y vectores propios: 1, 6a, 9a, 9c y 10b.}
\begin{preguntas}
%%%%%%%%%%%%%%%%%%%%%%%%%%%%%%%%%%%%%%%%%%%%%%%%%%%%%%%%%%%%%%%%%%%%%%%%%%%%%%%

%%%%%%%%%%%%%%%%%%%%%%%%%%%%%%%%%%%%%%%%%%%%%%%%%%%%%%%%%%%%%%%%%%%%%%%%%%%%%%%
\item Dados $a,b,c,d,e,f\in\R$, considere la matriz $A = \begin{pmatrix}1&a&b\\1&c&d\\1&e&f\end{pmatrix}$, una matriz simétrica. Determine los valores de $a,b,c,d,e,f$ tales que los vectores $(1,1,1)^{T}$ y $(1,0,-1)^{T}$ sean vectores propios de $A$, asociados a dos valores propios distintos.
    
\begin{respuesta}\hspace{0pt}

Por ser $A$ simétrica se tiene que $a = 1$, $b = 1$ y que $d = e$. Sean $\lambda_{1}$ y $\lambda_{2}$ los dos valores propios de $A$ que se consideran, entonces
\[
    \begin{pmatrix}1&1&1\\1&c&d\\1&d&f\end{pmatrix}\begin{pmatrix}1\\1\\1\end{pmatrix}=\lambda_{1}\begin{pmatrix}1\\1\\1\end{pmatrix}
    \texty
    \begin{pmatrix}1&1&1\\1&c&d\\1&d&f\end{pmatrix}\begin{pmatrix}1\\0\\-1\end{pmatrix}=\lambda_{2}\begin{pmatrix}1\\0\\-1\end{pmatrix}.
\]
A partir de lo cual se obtienen que $\lambda_{1}=3$, $\lambda_{2}=0$, $c=1$, $d=1$ y $f=1$.\qedhere
\end{respuesta}
%%%%%%%%%%%%%%%%%%%%%%%%%%%%%%%%%%%%%%%%%%%%%%%%%%%%%%%%%%%%%%%%%%%%%%%%%%%%%%%


\begin{respuesta}\hspace{0pt}

\begin{enumerate}
    \item Supongamos que existe otro valor propio $\lambda_{1}$ tal que $u$ es su vector propio asociado, entonces
    \[
        f(u) = \lambda_{1}u.
    \]
    Además, por hipótesis tenemos que
    \[
        f(u) = \lambda u.
    \]
    Por tanto, tenemos que $\lambda_{1}u=\lambda u$, lo que implica que $(\lambda_{1}-\lambda)u=0$, y dado que $u\neq0$, al ser un vector propio, concluimos que $\lambda_{1}=\lambda$. Por lo cual, el enunciado es verdadero.
    \item Sea $\alpha\in\R$, entonces
    \[
        f(\alpha u) = \alpha f(u) = \alpha (\lambda u) = \lambda (\alpha u).
    \]
    Por lo cual, $\alpha u$ es un vector propio de $f$ asociado al valor propio $\lambda$.
    \item Verifiquemos si $u+v$ es un vector propio de $f$; para esto debemos encontrar el correspondiente valor propio. En efecto,
    \[
        f(u + v) = f(u) + f(v) = \lambda u + \mu v.
    \]
    Si $\lambda = \mu$, tendríamos que $f(u + v) = f(u) + f(v) = \lambda(u+v)$, sin embargo, esto no es posible ya que se tiene que  $\lambda\neq\mu$, en efecto, notemos que la aplicación lineal
    \[
        \funcion{f}{\R^{2}}{\R^{2}}{(x,y)}{(x+2y,2x+y).}
    \]
    Cuya matriz asociada a la base canónica es $A = \begin{pmatrix}1&2\\2&1\end{pmatrix}$, cuyos valores propios son $-1$ y 3 con vectores propios asociados $(-1,1)^{T}$ y $(1,1)^{T}$, respectivamente. Además,
    \[
        f((-1,1)+(1,1))=f(0,2)=(4,2)=2(2,1).
    \]
    Por lo tanto, el enunciado es falso.\qedhere
\end{enumerate}
\end{respuesta}
%%%%%%%%%%%%%%%%%%%%%%%%%%%%%%%%%%%%%%%%%%%%%%%%%%%%%%%%%%%%%%%%%%%%%%%%%%%%%%%
%\item Sea $\func{f}{\R^{n}}{\R^{n}}$ una aplicación lineal de manera que $f \circ f = f$. Calcular los valores propios de $f$.
    
%\begin{respuesta}\hspace{0pt}

%Supongamos que $\lambda$ es un valor propio de $f$ y que $v$ es un vector propio asociado a este valor propio. Entonces, se tienen que
%\[
%    f(v) = \lambda v
%\]
%y que 
%\[
 %   (f \circ f)(v) = f(f(v)) = f(\lambda v) = \lambda f(v) = \lambda^{2} v.
%\]
%Por hipótesis tenemos que $f \circ f = f$, y, por tanto, 
%\[
 %   f(v) = (f \circ f)(v).
%\]
%Luego, 
%\[
%    \lambda v = \lambda^{2} v,
%\]
%lo que implica que $\lambda(1 - \lambda)v = 0$, y dado que $v\neq0$, al ser un vector propio, tenemos que $\lambda(1 - \lambda)=0$. Por lo tanto, los valores propios de $f$ son 0 y 1.\qedhere
%\end{respuesta}
%%%%%%%%%%%%%%%%%%%%%%%%%%%%%%%%%%%%%%%%%%%%%%%%%%%%%%%%%%%%%%%%%%%%%%%%%%%%%%%
\item Sea $f$ una aplicación lineal, hallar sus valores y vectores propios, donde:
\begin{enumerate}
    \item 
    \[
        \funcion{f}{\R^{2}}{\R^{2}}{(x,y)}{(x+y,2x-3y).}
    \]
    \item
    \[
        \funcion{f}{\Mat{2}{2}}{\Mat{2}{2}}{
        \begin{pmatrix}
        a&b\\
        c&d
        \end{pmatrix}
        }{\begin{pmatrix}
        a+b-d&b\\c&-a+c+d
        \end{pmatrix}
        .}
    \]
  
\end{enumerate}

\begin{respuesta}\hspace{0pt}

\begin{enumerate}
    \item Para el valor propio $\lambda_{1} = - \sqrt{6}-1$ su correspondiente vector propio es $\begin{pmatrix}-\sqrt{6}+2\\2\end{pmatrix}$; para el valor propio $\lambda_{2} = \sqrt{6}-1$ su correspondiente vector propio es $\begin{pmatrix}\sqrt{6}+2\\2\end{pmatrix}$.
    \item Sea $B$ la base canónica de $\Mat{2}{2}$, verifique que
    \[
        [f]_{B}^{B} = 
        \begin{pmatrix}
        1&1&0&-1\\
        0&1&0&0\\
        0&0&1&0\\
        -1&0&1&1
        \end{pmatrix}
    \]
    Verifique que la matriz $A$ cumple que para el valor propio $\lambda_{1} = 0$ su correspondiente vector propio es $\begin{pmatrix}1\\0\\0\\1\end{pmatrix}$; para el valor propio $\lambda_{2} = 1$ sus correspondientes vectores propios son $\begin{pmatrix}1\\0\\1\\0\end{pmatrix}$ y $\begin{pmatrix}0\\1\\0\\1\end{pmatrix}$; para el valor propio $\lambda_{3} = 2$ su correspondiente vector propio es $\begin{pmatrix}-1\\0\\0\\1\end{pmatrix}$.
    
    Ahora, para el valor propio $\lambda_{1} = 0$ se verifica que 
    \[
        [f(v_{1})]_{B} = \lambda_{1}[v_{1}]_{B}
    \]
    donde
    \[
        [v_{1}]_{B} = \begin{pmatrix}1\\0\\0\\1\end{pmatrix}
    \]
    de lo cual
    \[
        v_{1} = \begin{pmatrix}1&0\\0&1\end{pmatrix}
    \]
    
    Para el valor propio $\lambda_{2} = 1$ se verifica que 
    \[
        [f(v_{2})]_{B} = \lambda_{2}[v_{2}]_{B}
        \texty
        [f(v_{3})]_{B} = \lambda_{2}[v_{3}]_{B}
    \]
    donde
    \[
        [v_{2}]_{B} = \begin{pmatrix}1\\0\\1\\0\end{pmatrix}
        \texty
        [v_{3}]_{B} = \begin{pmatrix}0\\1\\0\\1\end{pmatrix}
    \]
    de lo cual
    \[
        v_{2} = \begin{pmatrix}1&0\\1&0\end{pmatrix}
        \texty
        v_{3} = \begin{pmatrix}0&1\\0&1\end{pmatrix}
    \]
    
    Finalmente, para el valor propio $\lambda_{3} = 2$ se verifica que 
    \[
        [f(v_{4})]_{B} = \lambda_{3}[v_{4}]_{B}
    \]
    donde
    \[
        [v_{4}]_{B} = \begin{pmatrix}-1\\0\\0\\1\end{pmatrix}
    \]
    de lo cual
    \[
        v_{4} = \begin{pmatrix}-1&0\\0&1\end{pmatrix}.
    \]
    \item Sea $B = \{1, t, t^{2}\}$ la base canónica de $R_{2}[t]$, notemos que la matriz de representación de $f$ respecto a la base $B$ tanto en la salida como en la llegada es
    \[
        A = [f]_{B}^{B} = \begin{pmatrix}1&-1&1\\-1&1&-1\\1&-1&1\end{pmatrix}.
    \]
    Verifique que la matriz $A$ cumple que para el valor propio $\lambda_{1} = 0$ sus correspondientes vectores propios son $\begin{pmatrix}1\\1\\0\end{pmatrix}$ y $\begin{pmatrix}-1\\0\\1\end{pmatrix}$; para el valor propio $\lambda_{2} = 3$ su correspondiente vector propio es $\begin{pmatrix}1\\-1\\1\end{pmatrix}$.
    
    Luego, para el valor propio $\lambda_{1} = 0$ se verifica que 
    \[
        [f(v_{1})]_{B} = \lambda_{1}[v_{1}]_{B}
        \texty
        [f(v_{2})]_{B} = \lambda_{1}[v_{2}]_{B}
    \]
    donde
    \[
        [v_{1}]_{B} = \begin{pmatrix}1\\1\\0\end{pmatrix}
        \texty
        [v_{2}]_{B} = \begin{pmatrix}-1\\0\\1\end{pmatrix}
    \]
    de lo cual
    \[
        v_{1} = 1 + t
        \texty
        v_{2} = -1 +  t^{2}.
    \]
    
    De forma similar, para el valor propio $\lambda_{2} = 3$ se verifica que 
    \[
        [f(v_{3})]_{B} = \lambda_{2}[v_{3}]_{B}
    \]
    donde
    \[
        [v_{3}]_{B} = \begin{pmatrix}1\\-1\\1\end{pmatrix}
    \]
    de lo cual
    \[
        v_{3} = 1 - t +  t^{2}.\qedhere
    \]
\end{enumerate}
\end{respuesta}
\begin{comment}
%%%%%%%%%%%%%%%%%%%%%%%%%%%%%%%%%%%%%%%%%%%%%%%%%%%%%%%%%%%%%%%%%%%%%%%%%%%%%%%
\item Sea $f$ una aplicación lineal, hallar sus valores y vectores propios, donde:
\begin{enumerate}
    \item 
    \[
        \funcion{f}{\R^{2}}{\R^{2}}{(x,y)}{(x+y,2x-3y).}
    \]
    \item
    \[
        \funcion{f}{\Mat{2}{2}}{\R_{3}[t]}{
        \begin{pmatrix}
        a&b\\
        c&d
        \end{pmatrix}
        }{(a+b-d) + bt + ct^{2} + (-a+c+d)t^{3}.}
    \]
    \item 
    \[
        \funcion{f}{\R_{2}[t]}{\R^{3}}{a+bt+ct^{2}}{(a-b+c,-a+b-c,a-b+c).}
    \]
\end{enumerate}

\begin{respuesta}\hspace{0pt}

Verificar que:
\begin{enumerate}
    \item Para el valor propio $\lambda_{1} = - \sqrt{6}-1$ su correspondiente vector propio es $\begin{pmatrix}-\sqrt{6}+2\\2\end{pmatrix}$; para el valor propio $\lambda_{2} = \sqrt{6}-1$ su correspondiente vector propio es $\begin{pmatrix}\sqrt{6}+2\\2\end{pmatrix}$.
    \item Para el valor propio $\lambda_{1} = 0$ su correspondiente vector propio es $\begin{pmatrix}1\\0\\0\\1\end{pmatrix}$; para el valor propio $\lambda_{2} = 1$ sus correspondientes vectores propios son $\begin{pmatrix}1\\0\\1\\0\end{pmatrix}$ y $\begin{pmatrix}0\\1\\0\\1\end{pmatrix}$; para el valor propio $\lambda_{3} = 2$ su correspondiente vector propio es $\begin{pmatrix}-1\\0\\0\\1\end{pmatrix}$.
    \item Para el valor propio $\lambda_{1} = 0$ sus correspondientes vectores propios son $\begin{pmatrix}1\\1\\0\end{pmatrix}$ y $\begin{pmatrix}-1\\0\\1\end{pmatrix}$; para el valor propio $\lambda_{2} = 3$ su correspondiente vector propio es $\begin{pmatrix}1\\-1\\1\end{pmatrix}$.\qedhere
\end{enumerate}
\end{respuesta}
\end{comment}
%%%%%%%%%%%%%%%%%%%%%%%%%%%%%%%%%%%%%%%%%%%%%%%%%%%%%%%%%%%%%%%%%%%%%%%%%%%%%%%
\item Dado $a\in\R\setminus\{0\}$, considere la matriz $A = \begin{pmatrix}1&a&a\\a&1&a\\a&a&1\end{pmatrix}$.
\begin{enumerate}
    \item Compruebe que $\begin{pmatrix}1\\1\\1\end{pmatrix}$ es un vector propio de $A$. Hallar el valor propio $\lambda$ asociado. Calcule una base de $E_{\lambda}$.
    \item Compruebe que $1-a$ es un valor propio de $A$ y calcule una base del espacio propio correspondiente.
    \item Compruebe que todo vector $x\in\R^{3}$ se puede escribir como $x=y+z$ donde $y\in E_{\lambda}$ y $z\in E_{1-a}$ y que $E_{\lambda} \cap E_{1-a} = \{0\}$.
\end{enumerate}
    
\begin{respuesta}\hspace{0pt}

\begin{enumerate}
    \item Verifique que el valor propio asociado es $\lambda=2a+1$ y que $E_{\lambda}=\gen(\{(1,1,1)\})$.
    \item Utilice el polinomio característico de $A - \mu I_{3}$ y compruebe que $E_{1-a}=\gen(\{(-1,1,0),(-1,0,1)\})$.
    \item Verifique que $E_{\lambda} + E_{1-a} = \R^{3}$ y concluya que $\dim(E_{\lambda} \cap E_{1-a}) = 0$.\qedhere
\end{enumerate}
\end{respuesta}
%%%%%%%%%%%%%%%%%%%%%%%%%%%%%%%%%%%%%%%%%%%%%%%%%%%%%%%%%%%%%%%%%%%%%%%%%%%%%%%
\item Sean $A,B\in\Mat{n}{n}$. Pruebe que si $\lambda$ es un valor propio de $AB$ también lo es de $BA$.

\emph{Sugerencia:} Distinga los casos $\lambda\neq0$ y $\lambda=0$. Al tratar este último, considere el polinomio característico de $AB$.
    
\begin{respuesta}\hspace{0pt}

Supongamos que $\lambda\neq0$, entonces, existe un vector $v$ no nulo tal que $(AB)v=\lambda v$. Ahora, sea $w = Bv$, tenemos que
\[
    (BA)w = B(Aw) = B(A(Bv)) = B((AB)v) = B(\lambda v) = \lambda Bv = \lambda w.
\]
Luego, $\lambda$ es un valor propio de $BA$.

Ahora, supongamos que $\lambda=0$, entonces
\[
    \det(AB - 0I_{n}) = \det(AB) = \det(BA) = \det(BA - 0I_{n}).
\]
Lo que implica que 0 también es valor propio de $BA$. Por lo tanto, todo valor propio de $AB$ también lo es de $BA$.\qedhere
\end{respuesta}
%%%%%%%%%%%%%%%%%%%%%%%%%%%%%%%%%%%%%%%%%%%%%%%%%%%%%%%%%%%%%%%%%%%%%%%%%%%%%%%
\item En cada caso encuentre los valores y vectores propios de la matriz dada.
\begin{multicols}{2}
\begin{enumerate}

\item $A_{1} = \begin{pmatrix}2&0&3\\0&1&0\\0&1&2\end{pmatrix}$

\begin{respuesta}\hspace{0pt} El polinomio característico es
\[
    p(\lambda) = - (\lambda - 1)(\lambda - 2)^{2}
\]
para cada $\lambda\in\R$. 

Para el valor propio 1 el vector propio asociado es
\[
    (-3a, - a, a)
\]
para cada $a\in\R\setminus\{0\}$.

Para el valor propio 2 el vector propio asociado es
\[
    (b, 0, 0)
\]
para cada $b\in\R\setminus\{0\}$.
\end{respuesta}
\item $A_{2} = \begin{pmatrix}0&0&0&1\\0&0&0&0\\0&0&0&0\\1&0&0&0\end{pmatrix}$
\end{enumerate}
\end{multicols}
%%%%%%%%%%%%%%%%%%%%%%%%%%%%%%%%%%%%%%%%%%%%%%%%%%%%%%%%%%%%%
\item En cada caso encuentre los valores y vectores propios de la matriz dada.
\begin{multicols}{2}
\begin{enumerate}
\item $A_{1} = \begin{pmatrix}1&-1\\2&4\end{pmatrix}$

\begin{respuesta}\hspace{0pt} El polinomio característico es
\[
    p(\lambda) = (\lambda - 2)(\lambda - 3)
\]
para cada $\lambda\in\R$. 

Para el valor propio 2 el vector propio asociado es
\[
    (-a, a)
\]
para cada $a\in\R\setminus\{0\}$.

Para el valor propio 3 el vector propio asociado es
\[
    (-b/2, b)
\]
para cada $b\in\R\setminus\{0\}$.
\end{respuesta}
\item $A_{2} = \begin{pmatrix}1&1&-2\\4&0&4\\1&-1&4\end{pmatrix}$

\end{enumerate}
\end{multicols}
%%%%%%%%%%%%%%%%%%%%%%%%%%%%%%%%%%%%%%%%%%%%%%%%%%%%%%%%%%%%%%


%%%%%%%%%%%%%%%%%%
\item Para cada matriz dada encuentre los valores  propios y correspondiente espacio propio asociado.
\begin{multicols}{2}
\begin{enumerate}

\item $A_{1} = \begin{pmatrix}0&1&1\\1&0&1\\1&1&0\end{pmatrix}$
\item $A_{2} = \begin{pmatrix}0&0&0&0\\0&0&1&0\\0&1&0&0\\0&0&0&0\end{pmatrix}$
\end{enumerate}
\end{multicols}
\item Para cada matriz dada encuentre los valores  propios y correspondiente espacio propio asociado.
\begin{multicols}{2}
\begin{enumerate}


\begin{respuesta}\hspace{0pt} El polinomio característico es
\[
    p(\lambda) = - \lambda^{2}(\lambda - 6)
\]
para cada $\lambda\in\R$. 

Para el valor propio 0 el espacio propio asociado es
\[
    E_{0} = \l\{(x,y,z)\in\R\,:\,x + y + z = 0\r\}.
\]

Para el valor propio 6 el vector propio asociado es
\[
    E_{6} = \l\{(x,y,z)\in\R\,:\,x = y, x = z \r\}.
\]
\end{respuesta}
\item $A_{1} = \begin{pmatrix}0&1&1\\1&0&1\\1&1&0\end{pmatrix}$
\item $A_{2} = \begin{pmatrix}0&0&0&0\\0&0&1&0\\0&1&0&0\\0&0&0&0\end{pmatrix}$
\end{enumerate}
\end{multicols}
%%%%%%%%%%%%%%%%%%%%%%%%%%%%%%%%%%%%%%%%%%%%%%%%%%%%%%%%%%%%%%%%%%%%%%%%%%%%%%%
\final
\end{preguntas}
\end{document}
